\documentclass{article}
\usepackage{iclr2027_conference,times}
\input{math_commands.tex}
\usepackage{hyperref}
\usepackage{url}
\usepackage{graphicx}
\usepackage{booktabs}
\usepackage{longtable}
\usepackage{amsmath,amssymb}

\title{How Accurate Is Neighbour-Fraction Regression for Global Effects Under Network Interference? Bias, Variance and Failure in a Controlled Simulation}
\author{Anonymous authors\\Paper under double-blind review}
%\iclrfinalcopy
\setcounter{topnumber}{3}\setcounter{totalnumber}{5}
\renewcommand\topfraction{0.9}\renewcommand\bottomfraction{0.8}\renewcommand\textfraction{0.1}\renewcommand\floatpagefraction{0.7}
\begin{document}
\maketitle
\begin{abstract}
We study the regression-based estimator of the global average treatment effect (GATE) under network interference in its common form: ordinary least squares of the outcome on a unit's own treatment and the treated fraction of its neighbours, extrapolated to the all-treated and all-control worlds.
Simulations with exactly known truth (806 planned cells over four synthetic and two Facebook graphs, five outcome mechanisms, four designs, four simulated worlds per cell; 172 cells added after the first results) separate its bias from its variance.
Under Bernoulli($p$) assignment and one-hop spillover, the population bias is $\gamma(s-1)$ for an exactly computable $s$; it matches 144 noiseless cells within Monte-Carlo error, does not shrink with $n$ ($+2.53\to+2.49$ from $n=10^3$ to $2\cdot10^4$), is proportional to the interference strength in 14 of 16 tested groups, and changes sign with $p$.
On Erd\H{o}s--R\'enyi graphs the standard deviation scales as $n^{-0.501\pm0.005}d^{0.495\pm0.006}$.
At the grid's $(n,d,\sigma)$ the estimator beats the difference in means once interference is moderate, but has 5.5--18.6 times its SD when there is none; the crossover strength scales as $n^{-0.49}d^{1.01}$ and the SD ratio as $d^{1.0}$.
Under nonlinear spillover bias dominates and richer fits (penalised OLS, tuned boosting) do not help; intervals that ignore dependence under-cover, and a cluster bootstrap does not reach 95\% either.
Heavy-tailed degrees raise neither SD nor bias in a matched test, and the Stanford3 case where graph clusters hurt replicates on a second seed.
Exposure-mapping baselines were estimated only on low-degree graphs and under cluster designs. Conclusions are limited to model-defined truth and these regimes.
\end{abstract}
\section{Introduction}
\label{sec:intro}
When users of a platform influence one another, the effect an A/B test is meant to inform is the \emph{global} average treatment effect (GATE):
the mean outcome if everybody were treated minus the mean outcome if nobody were. The usual difference in means estimates the effect of treating a unit
\emph{while its neighbours stay untreated}, and is biased for the GATE by the spillover. A popular remedy is the regression-based estimator of
\citet{chin2019} and \citet{hanugander2023}: regress outcomes on features of the assignment vector (own treatment, the treated fraction of one's neighbours, \dots),
predict the outcomes under all-treated and all-control assignment, and difference. Its guarantees are known for a correctly specified linear model
and in the large-sample limit; how accurate it is \emph{in practice}, what the error is made of, and when it fails is less clear.

We study these questions in a controlled simulation where the truth is known exactly: we define the potential-outcome functions, so the GATE of every simulated world is computed,
not estimated, and bias and variance can be separated and attributed. We vary interference strength and shape, graph family (synthetic and two
public Facebook graphs), size and degree, assignment design, noise, and the regression specification, always against a correctly specified control, and report
Monte-Carlo uncertainty on every number. Our contributions are:
\begin{itemize}\setlength\itemsep{0pt}
\item \textbf{A bias formula for the studied specification, and its test.} For the neighbour-fraction OLS$(z,f)$, one-hop spillover, Bernoulli($p$) assignment and the noiseless population limit, the bias is $\gamma(s-1)$ with $s$ a covariance ratio over the actual degree mixture (Eq.~\ref{eq:bias}). The principle behind it (a misspecified linear fit carries the local slope to the corners) is not new and is implicit in \citet{chin2019} and \citet{hanugander2023}; what we add is the closed form for finite and heterogeneous degrees, which predicts when the sign changes with $p$, and its test: it matches 144 simulated noiseless cells (all within 3 Monte-Carlo SE), does not shrink with $n$, and is proportional to $\gamma$ in a later grid (14 of 16 groups). A leverage heuristic $\mathrm{SD}\propto\sqrt{d/(n\,p(1-p))}$ is confirmed in its exponents ($-0.501\pm0.005$ in $n$, $0.495\pm0.006$ in $d$) on Erd\H{o}s--R\'enyi graphs.
\item \textbf{A map of when the estimator is accurate}, at the grids' sizes, degrees and noise levels. It beats the difference in means as soon as interference is moderate, but pays a 5--19$\times$ larger SD when there is none;
under nonlinear spillover its error is bias-dominated and sign-reversing in $p$; richer specifications (also penalised) and boosting mostly do not help; cluster
randomisation helps only on some graphs; and intervals that ignore dependence undercover, as does a cluster bootstrap.
\item \textbf{Negative and mixed results, reported as such}: the large-degree shortcut for the bias is a poor approximation on real graphs; the predicted penalty for heavy-tailed degrees was \emph{not} observed, also in a matched test; good graph clusters \emph{hurt} the regression estimator on the dense Stanford3 graph, which replicates on a second seed although our proposed mechanism does not hold there; Hájek exposure baselines are unestimable under Bernoulli assignment at the degrees of our main grids, and where estimable (low degree, cluster designs) the pre-stated rules for them gave ``unclear'' verdicts.
\end{itemize}
All results are generated by released code from fixed seeds (Appendix~\ref{app:repro}); conclusions are scoped to the simulated regimes: model-defined truth, one outcome family per mechanism, no real platform data.

\section{Related work}
\label{sec:related}
\textbf{Exposure mappings and design-based estimation.}
Design-based work defines the effect of an \emph{exposure} (a function of the unit's own and its neighbours' treatments) and estimates
contrasts of exposures by inverse-probability weighting, with variance bounds that hold for arbitrary interference once the exposure mapping is
correct \citep{aronowsamii2017}. The same ideas, with graph-cluster randomisation to make ``all neighbours treated'' likely, give Horvitz--Thompson
estimators whose variance grows with degree under independent assignment but can be controlled by clustering on restricted-growth graphs
\citep{ugander2013}; \citet{eckles2017} show that cluster randomisation and neighbour-aware analysis reduce the bias of GATE estimation in simulations,
most when the network is clustered and interference is strong, and \citet{karrer2021} describe an industrial deployment. \citet{chen2023} derive
bias and variance of an HT estimator under a one-hop model as a function of topology and the covariance of the assignment, and optimise the latter.
Our Hájek baselines follow this line; their failure when the exposure model is wrong is documented in these works and in \citet{chin2019}.

\textbf{Regression-based GATE estimation.}
\citet{chin2019} introduced the estimator we study: fit, by (separate-arm) least squares, a model of the outcome on features of the assignment of
\emph{other} units (e.g.\ the treated fraction and count of neighbours), predict outcomes at the all-treated and all-control feature values, and difference.
Under Bernoulli assignment, strict exogeneity and a correctly specified linear model it is unbiased with an exact variance expression whose
extrapolation term makes it several times more variable than the difference in means, and it comes with a distributional result that need not
be $\sqrt n$ under interference; flexible (additive-model) fits were reported to add little over OLS in one nonlinear simulation, and bootstrap
intervals rely on exogeneity.
\citet{hanugander2023} select graph- and treatment-based covariates from a larger generated set and give block-bootstrap intervals; their
guarantee is consistency (sparse linear truth, restricted-eigenvalue condition) rather than a finite-sample error expression, and they report that
simple fraction-only adjustment can be badly biased under misspecification and that naive bootstrap coverage can be well below nominal.
\citet{cortez2022} take a different route: under staggered rollout, GATE is estimable without network knowledge when interference is a low-degree
polynomial of the neighbours' treatments, which can be read as polynomial extrapolation of the same kind we study; its variance grows steeply in the degree of the polynomial.
We did not implement staggered rollout (Sec.~\ref{sec:limits}).

\textbf{What is new here.}
The cited works characterise the estimator under correct specification, prove consistency, or simulate particular networks and outcome models, and \citet{chin2019} and \citet{hanugander2023} already report that fraction-only adjustment can be biased under misspecification. We do not claim a new principle. We add (i)~a closed form of that bias for the neighbour-fraction OLS under one-hop spillover and Bernoulli assignment, evaluated on the actual degree mixture and tested against simulations with exactly known truth; (ii)~a controlled, quantitative map of error, coverage and baselines across graphs, mechanisms, designs and specifications, with bias and variance reported separately and a correctly specified control. Our positioning rests on the works cited here; we claim no priority beyond them.

\section{Setup: estimand, estimator, and where its error comes from}
\label{sec:setup}
\textbf{Estimand and model-defined truth.}
A graph $G$ on $n$ units has adjacency $A$, degrees $d_i$ and row-normalised $W=D^{-1}A$. An assignment $z\in\{0,1\}^n$ induces the
neighbour-treated fraction $f_i(z)=(Wz)_i$. The potential outcome $Y_i(z)$ may depend on the whole of $z$. The global average treatment effect is
$\tau_G=\frac1n\sum_i\,[Y_i(\mathbf 1)-Y_i(\mathbf 0)]$. Because we \emph{define} every outcome function, $\tau_G$ is computed exactly for each
simulated world (graph draw, baselines and noise); it is never estimated. All errors below are $\hat\tau-\tau_G$.

\textbf{Outcome mechanisms.}
The default is one-hop spillover $Y_i(z)=b_i+\tau z_i+\gamma\,h(f_i(z))+\varepsilon_i$ with direct effect $\tau=1$, strength $\gamma$, shape $h$ ($h(0)=0$, $h(1)=1$, so $\tau_G=\tau+\gamma$),
a baseline $b_i$ (graph-smooth, log-degree and i.i.d.\ parts) and noise $\varepsilon_i\sim N(0,1)$. Shapes are linear, concave (saturating), convex and sigmoid (threshold). We also use neighbour-count spillover, two-hop spillover, a degree-heterogeneous direct effect and
equilibrium propagation $Y=(I-\rho W)^{-1}(b+\tau z)+\varepsilon$ (Appendix~\ref{app:setup}).

\textbf{Estimators.}
The regression-based estimator fits OLS of outcomes on features $\phi_i(z)$ and predicts at the global treatment and control vectors:
$\hat\tau=\big(\tfrac1n\sum_i[\phi_i(\mathbf 1)-\phi_i(\mathbf 0)]\big)^{\!\top}\hat\beta$ \citep{chin2019,hanugander2023}.
Our main specification, \textbf{OLS$(z,f)$}, uses $\phi=(1,z_i,f_i)$; richer ones add the neighbour count, a two-hop fraction, log-degree
interactions, polynomials in $f$ (quad, cubic) or an 11-column set; a \textbf{correctly specified} (oracle) regression with the true
exposure columns is a control in every sweep, and gradient boosting is the flexible learner.
Baselines are the difference in means (DM) and Hájek exposure-mapping estimators (full-neighbourhood and $\kappa=0.75$) \citep{aronowsamii2017,chin2019}.

\textbf{Designs and graphs.}
Bernoulli$(p)$ assignment and graph-cluster randomisation with Louvain communities \citep{blondel2008} (Bernoulli$(0.5)$ per cluster), with random clusters of the same sizes and, on the SBM, planted blocks as controls.
Graphs: SBM \citep{holland1983}, Barab\'asi--Albert \citep{barabasi1999}, Watts--Strogatz \citep{watts1998} and Erd\H{o}s--R\'enyi (all $n=5000$, mean degree about 20; NetworkX \citep{hagberg2008}), and the real
Facebook100 Stanford3 \citep{traud2012,rossi2015} ($n=11{,}586$, mean degree 98) and SNAP ego-Facebook \citep{mcauley2012} ($n=4039$, mean degree 44); scaling grids vary $n\in[10^3,2\cdot10^4]$ and $d\in[5,80]$ (Appendix~\ref{app:setup}, \ref{app:repro}).

\textbf{Metrics.}
Per cell: bias, SD, RMSE, relative RMSE $=\mathrm{RMSE}/|\tau_G|$ and coverage of nominal 95\% Wald intervals (HC0, cluster-robust, graph-HAC over 1 or 2 hops), with
Monte-Carlo standard errors. The truth is exact, so $\mathrm{MSE}=\mathrm{bias}^2+\mathrm{SD}^2$ and the components are attributed separately.

\textbf{Where the error comes from.}
Three facts organise the study; the first two are made exact under stated conditions, the third is tested numerically.

\emph{(i) Bias is extrapolation of the local slope.} Under Bernoulli$(p)$ the own treatment $z_i$ is independent of $f_i$, which depends only on
neighbours. For noiseless one-hop outcomes the population least-squares coefficient on $f$ in OLS$(z,f)$ is
$s=\mathrm{Cov}(h(f),f)/\mathrm{Var}(f)$, the covariance being over the mixture of units with $f_i=\mathrm{Bin}(d_i,p)/d_i$. The estimator predicts $\tau+\gamma s$ whereas
the truth is $\tau+\gamma$, so
\begin{equation}\label{eq:bias}
  \mathrm{bias}_\infty \;=\; \gamma\,(s-1)\;\;\xrightarrow{\;d_i\to\infty\;}\;\;\gamma\,\big(h'(p)-1\big).
\end{equation}
The estimator transports the slope of the spillover curve at the operating point $p$ to the full range $[0,1]$. It is exact for $h$ linear (zero bias, for
any graph) and nonzero otherwise; it is a property of the model and the design and does not shrink with $n$ (Sec.~\ref{sec:res-theory}).
The derivation is in Appendix~\ref{app:theory}; for non-Bernoulli designs $z$ and $f$ are correlated and (\ref{eq:bias}) no longer applies.

\emph{(ii) Variance scales as $\sqrt{(1+d_h)/(n\,p(1-p))}$.} With errors of variance $\sigma_r^2$ treated as independent, $\mathrm{Var}(\hat\tau)\approx
\frac{\sigma_r^2}{n}\big[\frac1{p(1-p)}+\frac1{\mathrm{Var}(f)}\big]$ and $\mathrm{Var}(f)=p(1-p)\,\overline{d^{-1}}$, i.e.
$\mathrm{SD}\propto\sqrt{(1+d_h)/(n\,p(1-p))}$ with $d_h$ the harmonic mean degree. The exposure $f$ varies by only $O(1/\sqrt d)$ around $p$ whereas the
corners $f\in\{0,1\}$ lie $O(\sqrt d)$ standard deviations away, so the adjusted estimator pays a variance price that the difference in means does not.
The dependence among units that share neighbours changes the level, not the claimed scaling; both are tested.

\emph{(iii) Design and inference.} Cluster randomisation widens the support of $f$ (many units with $f\approx 0$ or $1$), which should reduce both
terms \emph{if} clusters follow the graph; and iid-style intervals ignore the dependence created by shared neighbours and by bias. We state the
predictions H1--H8 from the plan (Appendix~\ref{app:hyp}, which states when the numerical rules were fixed). Some of them (H1, H2, H3b, H6a, and the proportionality H3g of a later grid) are consequences of the model's algebra or checks of our code; we call them consistency checks and report them as passed or failed. The others are empirical, and are reported as supported, refuted or unclear.

\section{Experiments}
\label{sec:results}
Grids e1--e10 (Table~\ref{tab:grids}, Appendix) comprise 806 cells (about 58 CPU-hours); a cell has $4$ worlds $\times$ $100$ assignments ($150$ for the coverage grid; $2\times60$ and $2\times50$ for the noiseless grids), so most points
average 400 estimates. Error bars are $\pm1.96$ Monte-Carlo SE, in units of the direct effect $\tau=1$. Unless stated, the default scenario is SBM, $n=5000$, mean degree 20, Bernoulli(0.5), $\gamma=2$ ($\tau_G=3$).
Verdicts follow the decision rules of Appendix~\ref{app:hyp} (directions from the plan, thresholds fixed after the grids ran). Grids e11--e21 (816 cells, 1622 in all) were designed after the first results were seen (e11--e13 with decision rules, e14--e21 as estimates with uncertainty only). Between-world variance of OLS$(z,f)$ is small (intraclass correlation of the error: median $\le0.004$, maximum $0.035$; 20-world grid e18, Appendix~\ref{app:post2}); with four worlds, bias$\pm1.96\,$SE (larger of replicate- and world-level SE) covered the 20-world reference in $0.97$ of draws, the world-level SE alone in $0.85$.

\subsection{Accuracy versus the baselines (Q1)}
\label{sec:res-acc}
\begin{figure}[tbp]
\centering\includegraphics[width=0.55\linewidth]{fig1_strength.pdf}
\caption{Bias (top) and RMSE (bottom, log scale) against interference strength $\gamma$ (linear-in-means truth, Bernoulli(0.5)). The difference in means has bias exactly $-\gamma$; the regression estimators are unbiased, with a fixed variance price. Other graphs: Fig.~\ref{fig:strength-all}.}
\label{fig:strength}
\end{figure}
With linear spillover the neighbour-fraction regression is unbiased on all six graphs, whereas the difference in means has bias slope $-1.000$ to $-0.999$ per unit of $\gamma$ (Fig.~\ref{fig:strength}); it wins on RMSE at $\gamma=0$ on all 12 graph--design combinations, and the regression overtakes it from $\gamma=0.5$ (BA, ER, SBM, Stanford3) or $\gamma=1$ (WS, ego-Facebook) under Bernoulli assignment.
The price is variance: at $\gamma=0$, when both are unbiased, the SD of OLS$(z,f)$ is 5.5--18.6 times that of the difference in means (BA 5.5, ER 6.5, Stanford3 9.2, SBM 10.1, WS 12.3, ego-Facebook 18.6), and a correctly specified regression pays the same price (ratio 0.84--1.00 to OLS$(z,f)$): this is the cost of extrapolating to the corners, not of misspecification. Hierarchical-bootstrap intervals are narrow (BA 5.5 [4.8, 6.3], ego-Facebook 18.6 [16.9, 20.7]), so the range is variation over graphs. It is not a constant: on an SBM family with $n\in\{10^3,5\cdot10^3,2\cdot10^4\}$, mean degree 10--40 and noise $\sigma\in\{0.5,1,2\}$ (e17) the ratio is 4.0--26.5 and scales as $d^{1.0}\sigma^{-0.31}n^{0.01}$ ($R^2=0.99$), and the interference strength $\gamma^*$ at which the regression overtakes the difference in means scales as $n^{-0.49}d^{1.01}\sigma^{0.09}$ ($R^2=0.997$), from $0.11$ to $2.2$. In the default scenario with \emph{concave} spillover (Table~\ref{tab:est}) the difference in means has RMSE 2.00 (67\% of $\tau_G$), OLS$(z,f)$ 0.72 (24\%), the oracle 0.59 (20\%), and $\kappa$-Hájek 1.21; the neighbour-count regression has bias $+1.0$ (RMSE 1.05), and richer specifications trade bias for variance: the cubic fit removes it ($-0.05$) but has RMSE 1.33, the 11-column fit RMSE 3.01 (SD 2.95 against 0.49). Separate-arm fitting and penalisation do not change this on the 48 Bernoulli cells of e2 (e16, Appendix~\ref{app:post2}): the fit $(z,f,zf)$ has a median RMSE $1.03\times$ that of OLS$(z,f)$ and is better in 19\% of cells; ridge and lasso on the 11 columns have median ratios $6.3$ and $6.5$ and are never better; boosting tuned by cross-validation has median ratio $2.0$ at Bernoulli(0.5) and is better in 21\% of cells, mostly with sigmoid truth, since trees cannot extrapolate to $f\in\{0,1\}$.
\textbf{Exposure-mapping baselines (H6; post-freeze grid e12, Appendix~\ref{app:post}, Fig.~\ref{fig:hajekfig}).} Hájek estimators are mostly unestimable at the degrees above, so we added low-degree graphs where they are estimable: with a correct exposure mapping they are unbiased when both extreme exposures are common (Bernoulli(0.5) and clusters: 30 of 30 cells) but not at Bernoulli(0.8) (1 of 13); with a wrong mapping they are biased by the size of the mismatch; they beat OLS$(z,f)$ in RMSE in only about 30\% of cells. The pre-stated rules gave ``unclear'' for both parts (31 of 43 and 44 of 65 cells). Under Louvain-cluster designs at degree 20--98 (e19) the full-exposure Hájek estimator is nearly unbiased ($|\text{bias}|\le0.05$) but has RMSE $1.19$--$1.47$ times that of OLS$(z,f)$ on SBM and ego-Facebook, and is competitive only on Stanford3 with concave truth ($0.89$--$0.96$).
\begin{table}[tbp]
\caption{Estimators on the default scenario (SBM, Bernoulli(0.5), $\gamma=2$, $\tau_G=3$): bias (MC SE), SD, RMSE. Source: \texttt{results/summary/e2.csv}; generated by \texttt{figures/scripts/make\_tables.py}.}
\label{tab:est}
\centering\scriptsize\setlength\tabcolsep{3pt}
\input{tab2_estimators.tex}
\end{table}

\subsection{Where the bias comes from, and why it does not average out (Q2)}
\label{sec:res-theory}
\begin{figure}[tbp]
\centering\includegraphics[width=0.7\linewidth]{fig2_bias_theory.pdf}
\caption{(a) Noiseless simulated bias of OLS$(z,f)$ against the closed form of Eq.~(\ref{eq:bias}) in 144 cells (SBM, BA, Stanford3; concave/convex shapes, curvature 0.5--8, $p\in\{0.1,0.25,0.5,0.75\}$). (b,c) Bias against $p$: lines are the closed form, markers simulation.}
\label{fig:theory}
\end{figure}
\textbf{The bias is the extrapolation of the local slope (checks H1, H2).} With a correctly specified model and no noise, the bias is zero to machine precision (max $|\text{bias}|=2.4\times10^{-11}$ over 120 cells); in noisy cells only 0.6\% of 156 cells exceed 3 SE. For OLS$(z,f)$ under nonlinear one-hop spillover the exact expression (\ref{eq:bias}) reproduces all 144 noiseless cells (correlation $1.0000$, median $|\text{dev}|=0.0012$, max $0.041$, all within 3 SE; Fig.~\ref{fig:theory}, Table~\ref{tab:theory}) over a range of biases from $-2$ to $+5$ ($\tau_G=3$). The sign follows the curvature of $h$ at $p$: concave spillover at small $p$ is over-, at large $p$ under-estimated, and the sign flips with $p$. The large-degree shortcut $\gamma(h'(p)-1)$ is a much poorer description: it deviates by a median 0.09 and up to 3.3, on Stanford3 by 0.5--0.8 on average, because the low-degree units dominate the exposure distribution (the exact formula needs the degree distribution).
\textbf{It is not a sampling error (H3).} In SBM with concave spillover the bias of OLS$(z,f)$ is $+2.53$ at $n=10^3$ and $+2.49$ at $n=2\times10^4$ for $p=0.1$ ($-0.57\to-0.52$ at $p=0.5$; log--log slopes $-0.01$, $-0.03$), whereas the SD falls as $n^{-1/2}$ (Fig.~\ref{fig:bias-n}); more data tighten a wrong answer. The bias grows with the curvature (median Spearman 0.92 over 12 graph--design groups). The bias is proportional to the interference strength (post-freeze grid e11, Appendix~\ref{app:post}): 14 of 16 graph--shape--$p$ groups pass the rule (R$^2\ge0.99$, intercept within 3 SE; the two failures are ego-Facebook), with slopes per unit $\gamma$ for OLS$(z,f)$ from $-0.71$ to $+1.26$ set by shape and $p$.
\begin{figure}[tbp]
\centering\includegraphics[width=0.64\linewidth]{fig3_bias_vs_n.pdf}
\caption{Bias (symlog) and SD (log--log) against $n$, concave spillover, three designs. The oracle (green dashed) is unbiased; the bias of the misspecified regressions is flat in $n$ while the SD falls with slope $-1/2$ (dotted).}
\label{fig:bias-n}
\end{figure}

\textbf{Variance (H4).} On Erd\H{o}s--R\'enyi graphs under Bernoulli(0.5) the SD of OLS$(z,f)$ follows $\sqrt{d/n}$ with fitted exponents $-0.501\pm0.005$ in $n$ and $+0.495\pm0.006$ in $d$ over 25 cells (Fig.~\ref{fig:var}, Table~\ref{tab:scaling}); the correctly specified model and the model with a log-degree term have the same exponents, but the 11-column specification does not ($-0.61$, $+1.18$: its variance explodes with degree). The dependence on $p$ follows $1/\sqrt{p(1-p)}$ on SBM and Stanford3 (SD at $p=0.05$ relative to $p=0.5$: 2.32 and 2.21 against a predicted 2.29), but not on ego-Facebook (1.80; log--log slope $0.59$ against 1), whose hubs break the homogeneous-degree argument (decision rule: unclear). Under Louvain clusters the scaling differs ($n^{-0.55}$, $d^{+0.37}$). Raising the graph-smooth baseline component from 0 to 1 to 3 raises the SD of OLS$(z,f)$ from 0.16 to 0.49 to 1.37 at $\sigma=1$ (e5): neighbours' shared baselines are an additional variance source. The heuristic with the harmonic-mean degree $d_h$ (Sec.~\ref{sec:setup}) holds across graphs when the errors are independent (e21, $\gamma=0$, noise only: observed over predicted SD between $0.93$ and $1.04$ on all six graphs, while the arithmetic mean degree is off by a factor 2.8 between graphs), but not for the default outcome: with its graph-smooth and log-degree baselines the ratio ranges from $1.3$ (ER) to $5.1$ (ego-Facebook), and neither $d_h$ nor $d$ orders the six graphs (Spearman correlation $0.03$). The \emph{scaling} in $n$, $d$ and $p$ is thus supported for independent errors and homogeneous graphs; the \emph{level} on a given graph depends on the correlation of the baseline, which the formula does not contain.
\begin{figure}[tbp]
\centering\includegraphics[width=0.66\linewidth]{fig4_variance_scaling.pdf}
\caption{(a) SD against $n$ for each mean degree $d$ (ER, Bernoulli(0.5)); dotted: fitted $C\sqrt{d/n}$. (b) $\mathrm{SD}/\sqrt{d/n}$ is flat for OLS$(z,f)$, not for the 11-column fit. (c) SD relative to $p=0.5$ against $p$.}
\label{fig:var}
\end{figure}

\subsection{Design: when cluster randomisation helps (H5)}
\label{sec:res-design}
Louvain-cluster randomisation removes most of the difference-in-means bias in every matched cell ($|\text{bias}|$ smaller in 100\% of 36 matched resolutions $\times$ graph $\times$ shape cells), whereas random clusters of the same sizes leave it at the Bernoulli value ($-2.0$): clusters help through the graph, not through coarse assignment (Fig.~\ref{fig:cluster}, Appendix). For the regression estimator the picture depends on the graph. On SBM and ego-Facebook the Louvain design has lower RMSE (SBM, concave: bias $-0.07$ vs.\ $-0.48$, RMSE $0.37$ vs.\ $0.76$; ego-Facebook, linear: RMSE $0.43$ vs.\ $0.98$), but on the dense Stanford3 graph (mean degree 98) it is \emph{reversed} (RMSE $0.54$ vs.\ $0.26$, linear), and the ordering replicates on a second seed (Louvain worse in 10 of 12 matched cells in the first seed and 11 of 12 in the second; the 12-cell mean RMSEs flip, so the claim is about the cell-wise ordering). That good clusters leave little exposure variation to the regression holds for SBM and ego-Facebook (within-arm variance of $f$ 0.005 and 0.004--0.04, against 0.013 and 0.023 under Bernoulli) but \emph{not} for Stanford3, where it stays at the Bernoulli value (0.021--0.028 against 0.023); there the SD inflation goes with the few clusters and the cluster-level dependence of the treated share (Appendix~\ref{app:post2}). Overall Louvain beats random clusters in 72\% of 36 matched cells (decision rule $\ge75\%$: unclear).

\subsection{Inference (H7)}
\label{sec:res-cov}
\begin{figure}[tbp]
\centering\includegraphics[width=0.65\linewidth]{fig6_coverage.pdf}
\caption{Empirical coverage of nominal 95\% intervals for OLS$(z,f)$ (filled) and the correctly specified regression (open) by variance estimator (HC0; cluster-robust, CL, at analysis resolutions; graph-HAC at 1 and 2 hops), spillover shape and design.}
\label{fig:cov}
\end{figure}
Intervals that treat units as independent (HC0) undercover whenever neighbours share exposures. At the default analysis resolution, Bernoulli(0.5), and linear spillover, OLS$(z,f)$ has mean coverage 0.66 for HC0 (0.32--0.89 across the six graphs; lowest on ego-Facebook), 0.86 for cluster-robust errors, 0.79 for 1-hop graph-HAC and 0.91 for 2-hop graph-HAC (Fig.~\ref{fig:cov}). Under misspecification no interval works: with concave spillover the same estimator covers 0.43 (HC0), 0.65 (cluster), 0.56 (HAC-1) and 0.70 (HAC-2) on average, because the bias is not in the standard error; the correctly specified regression reaches 0.86 (cluster) and 0.92 (HAC-2), HC0 only 0.68. Binomial MC SEs of single coverage values are 0.01--0.02 (lower bounds). Cells with a single analysis cluster (BA and ER at analysis resolution 0.5) are degenerate and excluded; the 2-hop HAC of this grid is over five graphs.
A repeat on the same six graphs, with the Stanford3 2-hop HAC added, a Louvain-block bootstrap ($B=200$) and coverage of both the GATE and the expectation $E[\hat\tau]$ (which removes the bias; e15, mean over six graphs $\times$ linear/concave truth) gives, for GATE/$E[\hat\tau]$ under Bernoulli(0.5): HC0 $0.54/0.66$, 1-hop HAC $0.66/0.78$, cluster-robust $0.75/0.86$, block bootstrap $0.76/0.87$, 2-hop HAC $0.82/0.92$; and with Louvain clusters $0.36/0.37$, $0.56/0.57$, $0.80/0.82$, $0.83/0.85$, $0.73/0.74$ (Appendix~\ref{app:post2}). The gap between the targets is the bias; the bootstrap equals cluster-robust errors under Bernoulli assignment, is the best under cluster designs, and no interval reaches 95\% on average. We resampled analysis clusters; the bootstrap of \citet{hanugander2023} was not implemented.

\subsection{Failure map and ranking of factors (Q2, Q3)}
\label{sec:res-fail}
\begin{figure}[tbp]
\centering\includegraphics[width=0.72\linewidth]{fig7_failure_map_relrmse_sbm.pdf}
\caption{Relative RMSE (RMSE$/|\tau_G|$) for each estimator, spillover shape or mechanism and design on SBM (B.5/B.1: Bernoulli 0.5/0.1; Lv: Louvain clusters). Hatched: not estimable. Other graphs: Appendix~\ref{app:figs}.}
\label{fig:fail}
\end{figure}
\textbf{Shape and treatment probability.} Median $|\text{bias}|/\tau_G$ of OLS$(z,f)$ is 0.01 (linear), 0.14 (concave), 0.14 (convex) and 0.34 (sigmoid) over 24 cells each, and the squared bias exceeds half the MSE in 0, 33, 38 and 67\% of cells (Fig.~\ref{fig:fail}). At $p=0.1$ the regression family is much worse than at $0.5$ (on SBM, concave: relative RMSE 0.86 for OLS$(z,f)$, 3.8 for OLS quad and 27 for the 11-column set) and the bias changes sign (sigmoid: $+2.0$ at $p=0.5$, $-1.4$ to $-1.5$ at $p=0.1$). Over 96 graph--shape--design cells OLS$(z,f)$ beats the difference in means in 85\%, but a correctly specified model beats OLS$(z,f)$ in 85\%, while quadratic features (14\%) and gradient boosting (18\%) rarely do: flexibility costs variance and extrapolates badly (OLS quad bias $-11$ at $p=0.1$, $n=2\times10^4$).
\textbf{Other mechanisms.} With two-hop interference fitted by one-hop features OLS$(z,f)$ has bias $-0.4$ (WS) to $-1.0$ (SBM, BA, ER, Stanford3), yet \emph{adding} the two-hop exposure (correct model) has larger RMSE on every graph (SBM $2.75$ vs.\ $1.04$): the variance of the extra regressor exceeds the bias it removes. Near the singular end of equilibrium propagation ($\rho=0.9$; $\tau_G\approx10$, relative bias of OLS$(z,f)$ $-0.74$) no estimator is safe (mean RMSE: OLS$(z,f)$ 7.8, difference in means 8.9, boosting 8.7, 11-column 9.8); $\rho=0.3$ is benign.
\textbf{Graphs.} The predicted penalty for heavy-tailed degrees was not observed (H8): for linear truth OLS$(z,f)$ RMSE is 0.26 on BA and Stanford3, 0.27 ER, 0.47 SBM, 0.51 WS and 0.92 ego-Facebook, so the two heavy-tailed graphs are among the best, but graph family is confounded with $n$ and mean degree here. A matched test (e14: $n=5000$, mean degree 20, linear truth, $\gamma=2$) removes the confounding: the SD of OLS$(z,f)$ is 0.255 on BA ($m=10$), 0.248 on a configuration model with the same degree sequence, 0.276 on ER, 0.240 on a 20-regular graph and 0.144--0.175 on power laws (exponents 3.0, 2.5, 2.2; $d_h$ 3.6--4.5), so heavier tails lower it, in the direction of the $d_h$ heuristic, and there is no penalty; under concave truth they also lower the bias ($-0.12$ to $-0.14$ against $-0.45$ to $-0.52$).
\begin{figure}[tbp]
\centering\includegraphics[width=0.96\linewidth]{fig11_crossover_anova.pdf}
\caption{(a) Crossover strength $\gamma^*$ at which OLS$(z,f)$ overtakes the difference in means in RMSE (SBM, linear truth, Bernoulli(0.5); from the fitted difference-in-means bias slope and the OLS RMSE) against mean degree, lines by $n$, markers by $\sigma$; open markers: $\gamma^*>2$, extrapolated. (b) SD ratio SD(OLS$_f$)/SD(DM) at $\gamma=0$ on the same cells, with the six default graphs (95\% hierarchical-bootstrap intervals). Printed: exponents of the log--log fits over 27 cells ($\pm$ SE). (c) First-order and total-effect shares of the variance of $\log$ RMSE(OLS$_f$), SBM factorial of 144 cells (e20); whiskers: 95\% parametric bootstrap.}
\label{fig:cross}
\end{figure}
\textbf{Which factors matter (C3).} One-factor sweeps give max/min RMSE ratios of OLS$(z,f)$ of $4.6$ for sample size, $4.5$ for spillover shape, $3.9$ for mean degree, $3.6$ for $p$, $3.2$ for graph family, $2.0$ for design and $1.5$ for noise (Fig.~\ref{fig:rank}, Appendix~\ref{app:figs}), but these depend on the width of each range and are not a ranking. A full factorial on SBM (e20, 144 cells; $n$, $d$, $p$, $\sigma$, graph-smooth baseline, shape) decomposes the variance of log RMSE of OLS$(z,f)$ into first-order/total-effect shares: graph-smooth baseline $0.22/0.33$, spillover shape $0.20/0.36$, $p$ $0.16/0.20$, $n$ $0.13/0.17$, mean degree $0.09/0.12$, noise $0.01/0.04$ (Fig.~\ref{fig:cross}c). The outcome side (baseline structure and shape) thus matters most and noise least, in these ranges and on SBM; graph family and design were not crossed, and the one-factor ratios remain the only evidence for them.

\section{Discussion and limitations}
\label{sec:limits}
\textbf{What the evidence says.} The regression-based GATE estimator is accurate when its exposure model is right and the interference is not weak, and its failure is mostly a \emph{bias} problem that more data do not cure:
the bias is the slope of the spillover curve at the operating point carried to the corners, it is predicted by an exact formula in the one-hop Bernoulli case, and it changes sign with $p$.
The variance is a separate cost, paid even when the model is right (5.5--18.6 times the SD of the difference in means without interference) and growing with degree and with 1/$p(1-p)$.
The two causes respond to opposite remedies (a richer model removes bias and adds variance; more data remove variance only), so accuracy reports should give bias and SD separately and include a correctly specified control.
Standard errors that ignore dependence overstate precision, and no interval, including a cluster bootstrap, reaches nominal coverage on average once the bias is present.

\textbf{Limitations, most consequential first.} (i)~Truth is \emph{model-defined}: conclusions concern these mechanisms (one outcome family each, parameters at stated defaults), not real platforms. (ii)~Cells have 4 simulated worlds of 100 assignments; between-world variance of OLS$(z,f)$ is small (20-world check), but for the real graphs only baselines and assignments vary, and cell averages (e.g.\ coverage) describe the grid, not a population of problems. (iii)~Comparisons are at fixed $n$, degree and noise unless stated; the crossover and the SD ratio depend on them (Sec.~\ref{sec:res-acc}), and the factor decomposition covers SBM only. (iv)~The bias formula is derived and tested only for OLS$(z,f)$ with one-hop spillover, Bernoulli assignment and the noiseless population limit; the variance scaling holds for independent errors, its level on a given graph does not follow the formula. (v)~Richer fits were unpenalised, ridge/lasso-penalised or boosted with a small tuning grid; no other learner was tried. (vi)~Intervals are Wald intervals and an analysis-cluster bootstrap, not the covariate-selection bootstrap of \citet{hanugander2023}. (vii)~Hájek baselines were unestimable in most Bernoulli cells of the main grids and were studied on low-degree graphs ($n=3000$, degree 4--8) and under cluster designs; staggered rollout was not run. (viii)~For the real graphs, family is confounded with $n$ and degree; the Stanford3 reversal replicates on a second seed but its mechanism is not isolated. (ix)~First-grid thresholds were written after the grids ran (Appendix~\ref{app:hyp}); e14--e21 have none. (x)~The variance estimators have 16 unit tests; HAC-2 and the Hájek variance have no external check.

\subsection*{AI use statement}
This work was carried out by generative-AI agents (Anthropic Claude models) operating autonomously in a multi-agent research pipeline, with no human author writing text or code. The agents proposed and refined the hypotheses, designed the study, wrote and ran the simulation code, generated the synthetic outcome data, analysed the results, produced the figures and tables, searched and summarised the literature, formatted the references and wrote this paper; each part was done by a different agent role, and an independent review role checked deliverables against the task. Verification actually performed: (i) the bias formula was checked against 144 simulated noiseless cells and the correctly specified regression recovers the exact truth to $2.4\times10^{-11}$ in 120 noiseless cells; (ii) every number in the text is read from the committed result tables by a script (\texttt{paper/check\_numbers.py}) and the figures and tables are regenerated from the released code; (iii) bibliographic entries were checked against arXiv and Crossref metadata. What was \emph{not} done: no human reviewed the content line by line, and the two-hop graph-HAC and Hájek variance estimators were not checked against an independent implementation (the others were, in 16 unit tests). The system operator takes responsibility for the final content, including text, claims and artifacts produced with the aid of generative AI. See the ICLR 2027 AI Policy for Authors.

\subsection*{Ethics statement}
The study uses two public network datasets (Facebook100 Stanford3, distributed under CC BY-SA, and the SNAP ego-Facebook graph) and synthetic graphs; all outcomes are simulated from potential-outcome models that we define. There are no human-subject experiments, no new personal data, and no attempt to identify individuals; only anonymised edge lists are used. The methods concern estimating the effects of interventions on networks, which could be misapplied to platform experiments on users; our results mainly warn about overconfidence (biased point estimates and under-covering intervals), and we state the scope in which they hold.

\subsection*{Reproducibility statement}
All results come from the released code with fixed seeds (code, per-cell summaries and compact per-replicate data of all 21 grids, about 70 MB, are in the supplementary archive, so every table and figure can be regenerated without rerunning the simulations; the full raw per-replicate files are a separate release asset, with checksums in \texttt{results/MANIFEST.md}, and rerunning the simulations takes about 67 CPU-hours); Appendix~\ref{app:repro} gives the environment and library versions, data sources with SHA-256 hashes, seeds, the commands, measured runtimes and a map from every figure and table to its generating script. The exact grids and replicate counts are in Table~\ref{tab:grids}; hypotheses and decision rules are in Appendix~\ref{app:hyp}; the derivation of the bias formula is in Appendix~\ref{app:theory}.

\bibliography{refs}
\bibliographystyle{iclr2027_conference}
\appendix
\section{Setup details}
\label{app:setup}
\textbf{Outcome models.} All outcomes are $Y_i(z)=b_i+(\text{treatment part})+\varepsilon_i$ with a \emph{fixed} baseline per world, $b_i=u_i+0.5\,\tilde\ell_i+\eta_i$, where $u$ is a graph-smooth component (white noise passed through $W$ three times, rescaled to SD $\alpha_{\rm graph}$, default 1), $\tilde\ell_i$ the standardised log-degree, $\eta_i\sim N(0,0.5^2)$, and $\varepsilon_i\sim N(0,\sigma^2)$, $\sigma=1$. The mechanisms are:
\emph{shape}: $\tau z_i+\gamma h(f_i)$ with $h$ linear, concave $h(f)=\frac{1-e^{-kf}}{1-e^{-k}}$, convex $h(f)=\frac{e^{kf}-1}{e^{k}-1}$ ($k=3$ unless varied; $k\in[0.5,8]$ in the noiseless grid), or sigmoid (logistic with slope 10, rescaled to $h(0)=0,h(1)=1$);
\emph{count}: $\tau z_i+\gamma\,(Az)_i/\bar d$; \emph{twohop}: $\tau z_i+\gamma f_i+\gamma_2 f^{(2)}_i$ with $\gamma_2=0.5$ and $f^{(2)}$ the treated fraction among two-hop neighbours; \emph{hetero}: the direct effect $\tau(1+0.5\tilde\ell_i)$ with linear spillover; \emph{propagation}: $Y=(I-\rho W)^{-1}(b+\tau z)+\varepsilon$ for $\rho\in\{0.3,0.6,0.9\}$. Defaults are $\tau=1$, $\gamma=2$ (so $\tau_G=3$ for the shape mechanisms).
The GATE is the exact difference of the noise-free mean outcomes at $z=\mathbf 1$ and $z=\mathbf 0$ of the world.

\textbf{Estimators.} All OLS fits use the pseudo-inverse and predict at $z=\mathbf 1$ and $z=\mathbf 0$ with the same feature map. With $f_1$ the neighbour-treated fraction, $f_2$ the two-hop fraction, $c$ the neighbour count divided by the mean degree and $\ell$ the centred log-degree, the columns are: OLS$(z,f)$: $(1,z,f_1)$; OLS$_{\rm zf}$ adds $zf_1$; OLS$_{\rm cnt}$: $(1,z,c)$; OLS$_{\rm fd}$: $(1,z,f_1,\ell,z\ell,f_1\ell)$; OLS$_{f2}$: $(1,z,f_1,f_2)$; OLS quad: $(1,z,f_1,f_1^2,zf_1)$; OLS cubic adds $f_1^3$; OLS big: $(1,z,f_1,f_2,c,\ell,z\ell,f_1\ell,zf_1,f_1^2,f_2^2)$. ORACLE uses exactly the exposure columns of the mechanism ($(1,z,h(f_1))$ for shape, $(1,z,c)$ for count, $(1,z,f_1,f_2)$ for two-hop, $(1,z,z\ell,f_1)$ for heterogeneous direct effects); ORACLE$_D$ adds $\ell$, needed because baselines depend on degree; no oracle exists for propagation. GBM is scikit-learn histogram gradient boosting (150 iterations, depth 4, learning rate 0.1) on $(z,f_1,f_2,\log d,\text{count})$ with the same predict-at-corners contrast. HajekFull compares units that are treated with all neighbours treated against units that are control with all neighbours control, with exact design probabilities; HajekKappa ($\kappa=0.75$, Bernoulli only) uses ``treated with at least $\kappa$ of neighbours treated'' against ``control with at most $1-\kappa$ treated''. Both are undefined when no unit has the exposure; those cells are reported as unestimable (counts in the summaries).

\textbf{Variance estimators.} HC0 sandwich for the contrast; cluster-robust over fixed analysis clusters (Louvain partitions at resolution 0.5, 1, 2 or 5; 1 is the default); Conley-type graph-HAC with a uniform kernel over graph distance $\le1$ or $\le2$ (variance clipped at zero), computed only when $n\bar d^{2}<4\cdot10^{7}$, hence not for Stanford3's two-hop kernel.

\textbf{Grids and graphs.}
\begin{table}[h]
\caption{Grids: number of cells and replicates per cell (\texttt{results/tables/t13\_grid\_spec.csv}). $^\dagger$Designed after the first results were read (Appendix~\ref{app:post2}).}
\label{tab:grids}
\centering\scriptsize\setlength\tabcolsep{3pt}
\input{tab1_grids.tex}
\end{table}
\begin{table}[h]
\caption{Difference in means, OLS$(z,f)$ and the correctly specified regression by graph (concave spillover, $\gamma=2$, Bernoulli(0.5); \texttt{results/summary/e2.csv}).}
\label{tab:graphs}
\centering\scriptsize\setlength\tabcolsep{3pt}
\input{tab3_graphs.tex}
\end{table}
Stanford3: socfb-Stanford3, connected as downloaded ($n=11{,}586$, 568{,}309 edges, mean degree 98.1); ego-Facebook: largest connected component, $n=4039$, 88{,}234 edges, mean degree 43.7 (\texttt{results/tables/t13b\_real\_graphs.csv}). For real graphs only baselines, noise and assignments vary between worlds.

\begin{table}[h]
\caption{Closed form (\ref{eq:bias}) versus simulation, 144 noiseless cells; \texttt{results/summary/theory\_vs\_sim\_e9.csv}.}
\label{tab:theory}
\centering\scriptsize\setlength\tabcolsep{3pt}
\input{tab4_theory.tex}
\end{table}
\begin{table}[h]
\caption{Log--log fits of the SD on $n$ and mean degree $d$ (ER graphs, Bernoulli(0.5)); \texttt{results/summary/e3.csv}.}
\label{tab:scaling}
\centering\scriptsize\setlength\tabcolsep{3pt}
\input{tab5_scaling.tex}
\end{table}

\section{Derivations}
\label{app:theory}
\textbf{Bias formula (\ref{eq:bias}).} Take Bernoulli$(p)$ assignment, no self-loops, and noiseless outcomes $Y_i=b_i+\tau z_i+\gamma h(f_i)$. For unit $i$, $z_i$ is independent of $f_i=\mathrm{Bin}(d_i,p)/d_i$, and $\mathbb E f_i=p$ for all $i$. Pool units and assignments. Then $\mathrm{Cov}(z,f)=0$; $\mathrm{Cov}(b,\cdot)=0$ because $b_i$ is fixed and $f_i-p$ has mean zero given $i$; hence the least-squares coefficients on $(z,f)$ are $\tau$ for $z$ and $s=\mathrm{Cov}(h(f),f)/\mathrm{Var}(f)$ for $f$, with moments over the mixture of degrees. The predicted contrast is $\hat\tau\to\tau+s\gamma\,[f(\mathbf 1)-f(\mathbf 0)]=\tau+\gamma s$, since $f(\mathbf 1)-f(\mathbf 0)=1$, while $\tau_G=\tau+\gamma(h(1)-h(0))=\tau+\gamma$. This gives $\gamma(s-1)$. As $d_i\to\infty$, $f_i\to p$ and the delta method gives $\mathrm{Cov}(h(f),f)\approx h'(p)\mathrm{Var}(f)$, i.e.\ $s\to h'(p)$; for finite degrees $s$ is computed exactly from the binomial mixture (\texttt{code/scripts/theory\_bias.py}). For $h$ linear, $s=1$ and the bias is zero for any degree distribution. The sampling error of the OLS fit is an additional $O(n^{-1/2})$ term, which is why the formula is tested on noiseless cells with the finite-sample deviation reported alongside.

\textbf{Variance heuristic.} The contrast vector is $c=(0,1,1)$. With homoskedastic independent errors of variance $\sigma_r^2$ and $z$, $f$ uncorrelated, $\mathrm{Var}(\hat\tau)\approx\sigma_r^2\,c^\top(X^\top X)^{-1}c\approx\frac{\sigma_r^2}{n}\big[\frac{1}{p(1-p)}+\frac1{\mathrm{Var}(f)}\big]$, with $\mathrm{Var}(f)=p(1-p)\,\overline{1/d}$, so $\mathrm{Var}(\hat\tau)\approx\frac{\sigma_r^2}{n\,p(1-p)}(1+d_h)$, $d_h=1/\overline{1/d}$ the harmonic mean degree. We test the scaling in $n$, $d$, $p$ only, because neighbouring units share baselines and exposures, which violates the independence assumption and changes the level.

\section{Post-freeze grids e11--e13}
\label{app:post}
The three grids were run after the first results were seen, on a request from the planning step of the pipeline; the first-freeze summaries are unchanged. Rules are in \texttt{code/scripts/post\_freeze.py}, written (as reported by the simulation step; we could not verify the order from the commit history) before any cell of these grids was read; results are in \texttt{results/summary/e11.csv}, \texttt{e12.csv} and \texttt{results/tables/t14\_gamma\_dependence.csv}, \texttt{t15\_hajek.csv}.

\textbf{e11: bias against interference strength (H3, 64 cells, 400 estimates each).} SBM, BA, Stanford3 and ego-Facebook $\times$ concave and sigmoid spillover $\times$ $p\in\{0.1,0.5\}$ $\times$ $\gamma\in\{0.5,1,2,4\}$, noisy. The baseline and the noise do not depend on $\gamma$ and OLS is linear in the outcome, so the expected bias is exactly proportional to $\gamma$; the rule fits $\text{bias}=a+b\gamma$ per group and calls it proportional if $R^2\ge0.99$ and $|a|\le3\max(\mathrm{SE}_a,0.02)$. For OLS$(z,f)$, 14 of 16 groups are proportional (H3g supported). The slope per unit $\gamma$ is $-0.25$ (SBM, concave, $p=0.5$) and $+1.26$ ($p=0.1$), $+1.01$ and $-0.71$ for the sigmoid; BA $-0.23$, $+1.20$, $+0.89$, $-0.68$; Stanford3 $-0.08$, $+0.51$, $+0.29$, $-0.22$; ego-Facebook $-0.12$, $+0.88$, $+0.48$, $-0.40$ (slope SE 0.005--0.037). The sign and size are set by the shape and $p$ as in Sec.~\ref{sec:res-theory}; only the scale is $\gamma$. The two failing groups are both ego-Facebook (concave, $p=0.5$: $R^2=0.96$ with a small bias; sigmoid, $p=0.1$: intercept $+0.20\pm0.06$). The correctly specified regression with log-degree has $|\text{bias}|\le0.17$ at $\gamma=4$ in all groups, and the difference in means has slope $-1.00$. \emph{Negative result:} the quadratic regression OLS$(z,f,f^2)$ at $p=0.1$ has a $\gamma$-independent bias of $-9$ to $-10$ (SBM, BA), $-6$ (ego-Facebook) and $-1.9$ (Stanford3), an extrapolation instability (cf.\ e10), not a spillover effect.

\textbf{e12: Hájek baselines under correct and wrong exposure mappings (H6, 108 cells).} ER and WS (mean degree 4, 8) and BA ($m=2,4$), $n=3000$, Bernoulli(0.5), Bernoulli(0.8) and Louvain clusters; mechanisms: threshold (step) with $\kappa=1$ and $\kappa=0.75$ (the full-neighbourhood mapping is correct for both, the $\kappa=0.75$ mapping only for the second), linear and sigmoid one-hop shapes, two-hop, and propagation ($\rho=0.6$); noisy; $\tau_G=3$ for step and shape mechanisms. Both estimators are estimable in 89\% (full) and 92\% ($\kappa$) of cells ($n_{\text{rep}}\ge200$ of 400); all other cells are Bernoulli(0.8), where fully control-exposed units are rare. $\kappa$-Hájek is undefined under cluster designs. \emph{H6c (correct mapping):} unbiased ($|\text{bias}|\le\max(3\,\mathrm{SE},0.1)$) in 31 of 43 estimable cells (72\%; rule $\ge75\%$): unclear. By design (post hoc, same cells): Bernoulli(0.5) 18 of 18, clusters 12 of 12, Bernoulli(0.8) 1 of 13, where the ratio estimator is biased upward ($+0.5$ to $+0.7$ full, $+0.28$ $\kappa$) because of the few control-exposed units. \emph{H6d (wrong mapping):} biased ($>\max(3\,\mathrm{SE},0.2)$) in 44 of 65 cells (68\%): unclear; post hoc Bernoulli(0.5) 24 of 30, Bernoulli(0.8) 12 of 23, clusters 8 of 12. At Bernoulli(0.5) the mean bias is $-0.88$ (full on two-hop), $-0.65$ (full on propagation), $-0.57$ ($\kappa$ on linear) and $-1.38$ ($\kappa$ on the $\kappa=1$ step), but $-0.13$ for $\kappa$ on the sigmoid, because $f\ge0.75$ already lies on its saturated part. The full-neighbourhood estimator is not wrong on one-hop shapes (bias $0.03$ linear, $-0.01$ sigmoid): its full-exposure contrast equals the GATE for any shape. \emph{H6e (descriptive):} Hájek has lower RMSE than OLS$(z,f)$ in 32\% of 96 (full) and 30\% of 66 ($\kappa$) estimable cells (median RMSE ratio 1.10 and 1.28). The correctly specified regression with log-degree is unbiased ($|\text{bias}|\le0.04$) in the step, shape and two-hop mechanisms.

\begin{figure}[tbp]
\centering\includegraphics[width=\linewidth]{fig9_gamma_dependence.pdf}
\caption{Bias against interference strength $\gamma$ under concave and sigmoid spillover (post-freeze grid e11; marker = graph). The difference in means has bias $-\gamma$; OLS$(z,f)$ has bias proportional to $\gamma$ whose sign and size are set by shape and $p$ (14 of 16 groups proportional; both exceptions are ego-Facebook); the correctly specified regression is unbiased. OLS quad is omitted (Appendix~\ref{app:post}).}
\label{fig:gamma}
\end{figure}
\begin{figure}[tbp]
\centering\includegraphics[width=\linewidth]{fig10_hajek.pdf}
\caption{Bias relative to the GATE of the exposure-mapping (Hájek) estimators, the difference in means, OLS$(z,f)$ and the correctly specified regression on low-degree graphs (post-freeze grid e12), by mechanism and design (marker = graph). Filled: the exposure mapping or model is correct for the mechanism; open: wrong.}
\label{fig:hajekfig}
\end{figure}

\textbf{e13: second seed (217 cells).} The cells of the grids that carry H2, H4 and H5 were rerun with a new seed (\texttt{results/tables/t16\_second\_seed.csv}; the replication rule, at least 75\% of matched Stanford3 cells with Louvain RMSE above random-cluster RMSE, was fixed before any e13 cell was read, according to the simulation step). H5: Louvain RMSE exceeds random-cluster RMSE on Stanford3 in 10/12 matched cells in the first seed and 11/12 in the second, and in 0/12 on ego-Facebook in both (verdict ``replicated''); the cell-mean RMSEs flip between seeds (Louvain 0.92 vs 1.11 first seed, 0.99 vs 0.84 second), dominated by the 10-cluster resolution. H2: correlation of simulation and closed form $1.0000$, maximum deviation $0.058$, all 144 cells within 3 SE. H4: ER exponents $-0.506\pm0.008$ in $n$ and $+0.508\pm0.009$ in $d$.

\section{Post-freeze analyses e14--e21}
\label{app:post2}
These eight grids (427 cells) and their analyses were designed after the first results and after e1--e13 had been read, to answer questions raised on the first version. They have \emph{no} decision rule or threshold: each item is an estimate with its Monte-Carlo uncertainty. Summaries are \texttt{results/summary/e14.csv}--\texttt{e21.csv}, analysis tables \texttt{results/tables/t17}--\texttt{t25}, and per-replicate data for every grid are in \texttt{results/compact/} (Appendix~\ref{app:repro}).

\textbf{Variance law across graphs (e21 and \texttt{t17}).} Table~\ref{tab:vlaw}.  See Sec.~\ref{sec:res-theory}. With independent errors only ($\gamma=0$, no baseline), observed/predicted SD with $d_h$ is between 0.93 and 1.04 on all six graphs (max/min 1.12); with the arithmetic mean degree, max/min is 2.82 (Stanford3: 0.185 predicted, 0.064 observed). With the default outcome, observed/predicted is 1.29 (ER), 1.44 (BA), 2.40 (SBM), 2.60 (Stanford3), 2.59 (WS), 5.1 (ego-Facebook). The graph-smooth baseline multiplies the SD of OLS$(z,f)$ by 7.2 (ego-Facebook), 3.9 (WS), 3.7 (SBM), 3.6 (Stanford3), 2.1 (BA), 1.8 (ER). Stanford3 and ER have the same SD ($\approx0.26$) because a lower $d_h$ (10.8 against 19) is cancelled by a larger baseline effect.

\textbf{Stanford3 cluster-design mechanism (\texttt{t18}).} Table~\ref{tab:expo}.  Within-arm variance of $f$ under Louvain clusters of resolutions 1--20: Stanford3 0.021--0.028 (Bernoulli 0.023, random clusters of the same size 0.020--0.023) although $\mathrm{corr}(z,f)=0.70$--$0.88$; SBM 0.0046 (Bernoulli 0.013, $\mathrm{corr}=0.985$); ego-Facebook 0.0035--0.040 (Bernoulli 0.0225).

\textbf{Matched degree-distribution test (e14, \texttt{t23}).} Table~\ref{tab:heavy}.  28 cells, $n=5000$, mean degree 20, Bernoulli(0.5), 8 worlds $\times$ 50 assignments, linear and concave truth, $\gamma=2$. SD of OLS$(z,f)$ with linear truth and $\gamma=2$: BA ($m=10$) 0.255, configuration model on the BA degree sequence 0.248, ER 0.276, 20-regular 0.240, power laws with exponents 2.2, 2.5, 3.0 and degrees $\le500$: 0.144, 0.164, 0.175 (at $\gamma=0$: 0.259, 0.243, 0.263, 0.236, 0.147, 0.157, 0.163); $|\text{bias}|\le0.014$ for linear truth. Concave bias: BA $-0.452\pm0.012$, configuration model $-0.487\pm0.018$. The power laws also have a larger share of low-degree nodes ($d_h$ 3.6--4.5, clustering 0.05--0.08).

\textbf{Coverage (e15, \texttt{t19}, \texttt{t19b}).} Table~\ref{tab:cov2}.  Same 24 cells (6 graphs $\times$ linear/concave $\times$ 2 designs), $4\times150$ estimates, as e6 plus the block bootstrap and the Stanford3 2-hop HAC. Mean coverage of the GATE / of $E[\hat\tau]$ over the 12 cells of a design (Bernoulli(0.5); Louvain clusters): HC0 0.54/0.66; 0.36/0.37. HAC-1 0.66/0.78; 0.56/0.57. HAC-2 0.82/0.92; 0.73/0.74. Cluster-robust 0.75/0.86; 0.80/0.82. Block bootstrap (SD) 0.76/0.87; 0.83/0.85, (percentile) 0.75/0.86; 0.82/0.83. Minimum over cells of the GATE coverage: 0.31 (HC0, cluster-robust, bootstrap, Bernoulli).

On Stanford3 (concave/linear truth) the coverage of the GATE under Bernoulli is 0.93/0.95 (HAC-2), 0.84/0.87 (cluster-robust), 0.82/0.83 (HC0); under Louvain resolution 1, 0.57/0.59 (HAC-2), 0.73/0.76 (cluster-robust) and 0.84--0.92 (bootstrap).

\textbf{Estimator scope (e16, \texttt{t25}).} Table~\ref{tab:scope}.  The 48 Bernoulli cells of e2 rerun with the same seeds (OLS$(z,f)$, OLS$(z,f,zf)$ and OLS$_\text{big}$ reproduce e2 to the last digit). RMSE relative to OLS$(z,f)$, median over cells: OLS$(z,f,zf)$ 1.03 (1.00 at $p=0.5$, 1.09--1.72 at $p=0.1$); ridge 6.3 and lasso 6.5, against 6.6 for unpenalised OLS$_\text{big}$, and never below 1; boosting with depth $\{2,3,6\}$, learning rate $\{0.03,0.1,0.3\}$ and $\{100,300\}$ iterations chosen by cross-validation on the first replicate of each world: 2.03 at Bernoulli(0.5) (24 cells), against 1.71 for the untuned boosting over 48 cells; it wins only where OLS$(z,f)$ is badly biased (concave at $p=0.1$: 0.51; sigmoid: 0.40--0.59). Median $|\text{bias}|$ of tuned boosting 0.44 against 0.82 untuned.

\textbf{Crossover and variance ratio (e17, \texttt{t21}).} 135 cells: SBM, Bernoulli(0.5), linear truth, $n\in\{10^3,5\cdot10^3,2\cdot10^4\}$, $d\in\{10,20,40\}$, $\sigma\in\{0.5,1,2\}$, $\gamma\in\{0,0.25,0.5,1,2\}$. The bias slope of the difference in means is $-1.006$ to $-0.995$ in all 27 groups. $\gamma^*$ is the model-based crossover of RMSE (validated against interpolation of observed RMSEs); its values above the $\gamma$ grid (maximum 2.2 at $n=10^3$, $d=40$, $\sigma=2$) are extrapolations. Log--log fits: $\gamma^*\propto n^{-0.49\pm0.01}d^{1.01\pm0.02}\sigma^{0.09\pm0.02}$ ($R^2=0.997$), SD ratio at $\gamma=0$ $\propto n^{0.01\pm0.01}d^{1.01\pm0.02}\sigma^{-0.31\pm0.02}$ ($R^2=0.989$; 4.0 to 26.5, median 11.1).

\textbf{Variance decomposition (e20, \texttt{t22}).} SBM full factorial: $n\{10^3,5\cdot10^3,2\cdot10^4\}\times d\{10,20,40\}\times p\{0.1,0.5\}\times\sigma\{0.5,2\}\times$ graph-smooth baseline SD $\{0,2\}\times$ truth $\{$linear, concave$\}$, $\gamma=2$. First-order/total-effect shares (Sobol indices of cell means; MC uncertainty $\pm0.01$) for log RMSE of OLS$(z,f)$: baseline $0.22/0.33$, shape $0.20/0.36$, $p$ $0.16/0.20$, $n$ $0.13/0.17$, $d$ $0.09/0.12$, $\sigma$ $0.01/0.04$; the largest interaction is baseline $\times$ shape (0.06). For the SD: baseline $0.46/0.51$, $n$ $0.27/0.27$, $d$ $0.14/0.16$, $p$ $0.05/0.05$, $\sigma$ $0.03/0.06$, shape $0.00/0.00$. Two levels for four of the six factors; the shares describe the chosen levels.

\textbf{Uncertainty with four worlds (e18, \texttt{t20}).} Table~\ref{tab:worlds}.  12 cells (six graphs $\times$ linear/concave, Bernoulli(0.5)) with 20 worlds $\times$ 25 assignments (500 estimates). Intraclass correlation of the OLS$(z,f)$ error across worlds: median 0.000--0.004, maximum 0.035 (e18, concave) and 0.016 (e2); for the correctly specified regression up to 0.165 (e18) and 0.105 (e2, concave); the difference in means $\approx0$. Using subsets of $W$ worlds against the remaining worlds as reference, the interval bias$\pm1.96\,$SE with the aggregation rule (larger of replicate- and world-level SE) covers the reference in $0.968$ of draws at $W=4$, the world-level SE alone in $0.850$ ($0.940$ with a $t_3$ quantile; its 10th--90th percentile range is 0.014--0.079 around a median 0.031), the replicate-level SE alone in $0.958$; at $W=12$ the world-level SE covers $0.921$ ($0.946$ with $t$). This analysis was run for OLS$(z,f)$ only; with larger intraclass correlation (correctly specified model, concave) the world-level term matters more. Hierarchical-bootstrap 95\% intervals for the SD ratio of OLS$(z,f)$ to the difference in means at $\gamma=0$ (Bernoulli(0.5), $B=400$): SBM 10.1 [9.3, 11.1], BA 5.5 [4.8, 6.3], WS 12.3 [11.3, 13.6], ER 6.5 [5.8, 7.4], Stanford3 9.2 [7.9, 10.5], ego-Facebook 18.6 [16.9, 20.7].

\textbf{Design-based baselines under cluster designs at higher degree (e19, \texttt{t24}).} Louvain resolutions 1 and 5; SBM ($d=20$), ego-Facebook (44), Stanford3 (98); linear/concave, $\gamma=2$, 400 estimates. The full-exposure Hájek estimator is estimable in all but 2--6 of 400 replicates (ego-Facebook), with $|\text{bias}|\le0.05$ and RMSE relative to OLS$(z,f)$ of 1.19--1.25 (SBM), 1.32--1.47 (ego-Facebook), 0.89--0.96 (Stanford3, concave) and 1.03/2.01 (Stanford3, linear, resolutions 1/5); the correctly specified regression has the lowest RMSE in 11 of 12 cells (the difference in means in the twelfth). Fractional-exposure Hájek estimators ($\kappa=0.75,0.9$) are biased ($-0.04$ to $-0.77$) because they estimate a contrast at exposure $\kappa$, not 1. The first-version statement that design-based baselines are impractical at $d\ge20$ holds for Bernoulli designs only.

\begin{table}[h]
\caption{Coverage of nominal 95\% intervals of OLS$(z,f)$ for the GATE and for $E[\hat\tau]$ by interval type (e15): mean over 12 cells (six graphs $\times$ linear/concave, $\gamma=2$) with a bootstrap 95\% interval over cells, and min--max. Source: \texttt{t19b\_coverage\_by\_se\_type.csv}.}
\label{tab:cov2}
\centering\scriptsize\setlength\tabcolsep{3pt}
\resizebox{\linewidth}{!}{\input{tab6_coverage_se_types.tex}}
\end{table}
\begin{table}[h]
\caption{Observed SD of OLS$(z,f)$ against the variance law with harmonic-mean degree $d_h$ and mean degree $d$; $r$ = observed/predicted (e1, e3, e21). Source: \texttt{t17*.csv}.}
\label{tab:vlaw}
\centering\scriptsize\setlength\tabcolsep{3pt}
\resizebox{\linewidth}{!}{\input{tab7_variance_law.tex}}
\end{table}
\begin{table}[h]
\caption{Matched degree-distribution test (e14; $n=5000$, mean degree 19--20, Bernoulli(0.5)): SD and its ratio to the $d_h$ law at $\gamma=0$; bias (MC SE), SD and RMSE at $\gamma=2$. Source: \texttt{t23\_heavy\_tail.csv}.}
\label{tab:heavy}
\centering\scriptsize\setlength\tabcolsep{3pt}
\resizebox{\linewidth}{!}{\input{tab8_heavy_tail.tex}}
\end{table}
\begin{table}[h]
\caption{Estimator scope (e16): median RMSE relative to OLS$(z,f)$ and, after the slash, share of cells in which the estimator beats it, by spillover shape and $p$ (six graphs per entry); last column over all cells (number of cells). Ratios above 1 favour OLS$(z,f)$; ``--'': not run. Source: \texttt{t25\_estimator\_scope.csv}.}
\label{tab:scope}
\centering\scriptsize\setlength\tabcolsep{3pt}
\resizebox{\linewidth}{!}{\input{tab9_estimator_scope.tex}}
\end{table}
\begin{table}[h]
\caption{Variation of the neighbour-treated fraction $f$ under Bernoulli(0.5), Louvain-cluster and size-matched random-cluster designs (res.: resolution): variance of $f$, within-arm variance, $\mathrm{corr}(z,f)$ and SD of OLS$(z,f)$ in e4. Source: \texttt{t18\_exposure\_variation.csv}.}
\label{tab:expo}
\centering\scriptsize\setlength\tabcolsep{3pt}
\resizebox{\linewidth}{!}{\input{tab10_exposure_variation.tex}}
\end{table}
\begin{table}[h]
\caption{(a) Subsets of $W$ of the 20 worlds of e18 (12 cells): median SE of the bias by the aggregation rule, world-level only and replicate-level only, and how often bias $\pm1.96$ SE (or $\pm t$ SE) contains the 20-world mean. (b) Intraclass correlation of the error between worlds and design effect (DE) of the replicate-level SE, median and maximum over 12 cells per row. Sources: \texttt{t20b\_worlds\_needed.csv}, \texttt{t20\_variance\_components.csv}.}
\label{tab:worlds}
\centering\scriptsize\setlength\tabcolsep{3pt}
\resizebox{\linewidth}{!}{\input{tab11_worlds_icc.tex}}
\end{table}

\section{Hypotheses, decision rules and verdicts}
\label{app:hyp}
\textbf{When the rules were fixed.} The directional predictions H1--H8 were written in the study plan before any simulation was run. The numerical thresholds of the decision rules were written into \texttt{code/scripts/findings.py} \emph{after} the grids e1--e10 had completed, and were not changed to resolve ``unclear'' verdicts; the verdicts are therefore not pre-registered tests in the strict sense. The rules for the post-freeze grids e11 and e12 were written in \texttt{code/scripts/post\_freeze.py} before any of their cells was read (a statement we could not verify from the commit history, which does not order the two); those grids were nevertheless run after the first results were seen, at the request of a reviewing step of the pipeline, and are labelled as additions. Rules are those in \texttt{results/hypotheses.csv} and \texttt{results/hypotheses\_postfreeze.csv}. Part~A of the table lists consistency checks (identities of the model or checks of the code, which are expected to pass; ``passes'' says our implementation agrees with the algebra, not that the world behaves a certain way); parts~B and~C list empirical hypotheses, with verdicts ``supported'', ``refuted'' or ``unclear'' (the rule neither confirms nor refutes). Descriptive tests have no threshold.
{\scriptsize\begin{longtable}{@{}p{0.11\linewidth}p{0.33\linewidth}p{0.4\linewidth}p{0.08\linewidth}@{}}
\toprule ID & test & decision rule & verdict \\
\endhead
\multicolumn{4}{@{}l}{\emph{A. Consistency checks: identities of the model or of the code (a failure would indicate an error, a pass is not evidence about the world)}} \\ \midrule
H1 & noiseless ORACLE bias, e8 (all graphs/designs/shapes run) & supported if max$|$bias$|$$<$1e-8 (machine precision); refuted if any $|$bias$|$$>$5 MC SE & passes \\
H1 & noisy ORACLE\_D bias in units of its MC SE, e1,e2,e3,e5 & supported if share $|$z$|$$>$3 below 0.05 (expect $\sim$0.003 under H0) & passes \\
H2 & e9 noiseless OLS\_f bias vs exact-binomial closed form (theory\_bias.py) & supported if corr$>$0.99 and median$|$dev$|$$<$0.01 (bias scale 0.1-5); report share within 3 MC SE & passes \\
H3b & e9 noiseless $|$bias$|$ vs curvature (gamma fixed) & supported if median rho$>$0.8; gamma-dependence is NOT tested here (e9 fixes gamma) - see note & passes \\
H6a & e1 DM bias slope in gamma (truth: GATE = tau + gamma; DM sees tau only under Bernoulli) & supported if Bernoulli slopes within 0.05 of -1 (proportional to gamma) & passes \\
H3g & e11 OLS\_f: bias vs gamma in {0.5,1,2,4}, per graph x shape x design (post-freeze) & per group weighted line bias = a + b*gamma; proportional if R2 $\ge$ 0.99 and $|$a$|$ $\le$ 3 max(SE\_a, 0.02); supported if $\ge$ 75\% of groups with $|$bias(4)$|$$>$0.2 & passes \\
\midrule \multicolumn{4}{@{}l}{\emph{B. Empirical hypotheses (first freeze)}} \\ \midrule
H3a & e10 concave spillover: OLS\_f bias vs n (SBM, 1k-20k, Bernoulli) & supported if $|$bias$|$ at the largest n is $\ge$70\% of that at n=1000 for both p (sampling error would fall like n\^{}-1/2, to $\sim$22\%) & supported \\
H4a & e3 ER Bernoulli(0.5), OLS\_f: log SD $\sim$ log n + log d & supported if both exponents within 0.05 of prediction & supported \\
H4b & e3 SBM + Louvain-cluster design, OLS\_f: log SD $\sim$ log n + log d & descriptive (the prediction is stated for Bernoulli designs) & unclear \\
H4c & e4 SD vs p, OLS\_f, linear truth & supported if slope within [0.7,1.3] (pred 1) for every graph & unclear \\
H5a & e4 matched-resolution comparison of Louvain vs random clusters (OLS\_f RMSE; DM $|$bias$|$) & supported if good clusters have lower RMSE in $\ge$75\% of matched cells and random clusters show no gain over Bernoulli & unclear \\
H5b & e4 planted-block quality (within-block edge fraction 0.8/0.9/0.97), linear truth & descriptive: monotone improvement in RMSE with quality would support H5 & unclear \\
H6b & e1 DM wins RMSE at gamma=0, loses at larger gamma & supported if DM wins at gamma=0 in all combinations and OLS\_f wins by gamma$\le$4 & supported \\
H6c & Hajek / exposure-mapping baselines estimable? & reported as unestimable when n\_rep$<$10 (full-neighbourhood exposure units absent at this degree) & unclear \\
H7a & e6 mean 95\% CI coverage by SE type (degenerate cell e6 BA acl0.5 excluded: one analysis cluster) & supported if HC0 coverage $<$0.85 for correctly specified/linear cells and cluster-robust/HAC raise it; coverage loses under bias (concave OLS\_f $<$ 0.85) & supported \\
H8a & e2 graph family: OLS\_f RMSE (Bernoulli 0.5, linear truth) & supported if both heavy-tailed graphs (BA, Stanford3) have RMSE above the median of the six graphs; refuted if both are below it & refuted \\
H8b & e7 two-hop interference fitted with 1-hop features (OLS\_f) vs the 2-hop-aware ORACLE\_D, Bernoulli 0.5 & supported if OLS\_f $|$bias$|$$>$0.3 on all graphs & supported \\
C2a & e1 gamma=0, Bernoulli(0.5): SD of OLS\_f over SD of DM (no interference, both unbiased) & supported if the ratio exceeds 3 on every graph & supported \\
C2b & e2 OLS\_f: bias vs variance by spillover shape (gamma=2, 6 graphs x 4 designs) & bias dominates for sigmoid, not for linear (bias-dominated share 0 for linear) & supported \\
C2c & e2 RMSE comparisons over 96 graph x shape x design cells & descriptive (ranking of estimators by RMSE) & unclear \\
\midrule \multicolumn{4}{@{}l}{\emph{C. Empirical, post-freeze additions (run after the first results were seen)}} \\ \midrule
H6c & e12 Hajek bias under a CORRECT exposure mapping (step mechanism) & supported (unbiased) if $\ge$ 75\% of estimable (n\_rep $\ge$ 200) cells & unclear \\
H6d & e12 Hajek bias under a WRONG exposure mapping (HajekFull on two-hop/propagation; HajekKappa(0.75) on linear/sigmoid shape and on step kappa=1) & supported (biased) if $\ge$ 75\% of estimable cells & unclear \\
H6c-posthoc & H6c stratified by design (post-hoc reading of the same cells; the verdict above is the pre-stated rule) & descriptive & reported \\
H6d-posthoc & H6d stratified by design (post-hoc) & descriptive & reported \\
H6e & e12 RMSE of Hajek vs OLS\_f over estimable cells & descriptive (no verdict threshold) & reported \\
\bottomrule
\end{longtable}
}

\section{Reproduction and provenance}
\label{app:repro}
\textbf{Figure and table map} (paper number: file in \texttt{figures/out/}). Figs.~\ref{fig:strength}: \texttt{fig1\_strength}; \ref{fig:theory}: \texttt{fig2\_bias\_theory}; \ref{fig:bias-n}: \texttt{fig3\_bias\_vs\_n}; \ref{fig:var}: \texttt{fig4\_variance\_scaling}; \ref{fig:cov}: \texttt{fig6\_coverage}; \ref{fig:fail}: \texttt{fig7\_failure\_map\_relrmse\_sbm}; \ref{fig:rank}: \texttt{fig8\_factor\_ranking}; \ref{fig:cross}: \texttt{fig11\_crossover\_anova}; \ref{fig:gamma}: \texttt{fig9\_gamma\_dependence}; \ref{fig:hajekfig}: \texttt{fig10\_hajek}; \ref{fig:strength-all}: \texttt{fig1\_strength\_all}; \ref{fig:cluster}: \texttt{fig5\_cluster\_design}; Figs.~\ref{fig:fail-ba}--\ref{fig:fail-egofacebook}: \texttt{fig7\_failure\_map\_relrmse\_}\{\texttt{ba, er, ws, stanford3, egofacebook}\}. Tables~\ref{tab:est}: \texttt{tab2\_estimators}; \ref{tab:grids}: \texttt{tab1\_grids}; \ref{tab:graphs}: \texttt{tab3\_graphs}; \ref{tab:theory}: \texttt{tab4\_theory}; \ref{tab:scaling}: \texttt{tab5\_scaling}; \ref{tab:cov2}: \texttt{tab6\_coverage\_se\_types}; \ref{tab:vlaw}: \texttt{tab7\_variance\_law}; \ref{tab:heavy}: \texttt{tab8\_heavy\_tail}; \ref{tab:scope}: \texttt{tab9\_estimator\_scope}; \ref{tab:expo}: \texttt{tab10\_exposure\_variation}; \ref{tab:worlds}: \texttt{tab11\_worlds\_icc}. Each is produced by the script of the same name (tables: \texttt{make\_tables.py}) in \texttt{figures/scripts/} from \texttt{results/summary/}; \texttt{figures/MANIFEST.md} lists inputs. \textbf{Environment.} Python 3.11.2; numpy 2.4.6, scipy 1.17.1, pandas 3.0.6, scikit-learn 1.9.1, networkx 3.6.1 (\texttt{code/requirements.lock}); CPU only. \textbf{Data.} socfb-Stanford3 zip SHA-256 \texttt{b1221e41\dots8720d} and SNAP ego-Facebook SHA-256 \texttt{125e84db\dots035d7}; outcomes are simulated, no outcome data exist. \textbf{Per-replicate data.} \texttt{results/compact/} holds one row per cell, world, replicate and estimator for all 21 grids (about 70 MB); \texttt{aggregate.py --compact} rebuilds every per-cell summary from it, equal to the shipped summaries to a relative $3.5\cdot10^{-7}$, so all tables and figures can be regenerated without rerunning the simulations. \textbf{Seeds.} Each cell has a fixed integer seed; world $w$ uses \texttt{default\_rng(seed*1000+w)}. \textbf{Commands and runtime.} \texttt{REPRODUCE.md} lists the commands (run experiment, aggregate, tables, findings, factor ranking, figures); grids e1--e10 take about 58 CPU-hours (e9 1.4 h, e2 15 h, e7 11 h) and the post-freeze grids e11--e21 about 8 more (e13 3.3 h, e16 2.5 h, the others under 0.7 h each; \texttt{results/RUNTIME.csv}). The unit tests (16, 9 s) run with \texttt{pytest code/tests}. \textbf{Provenance of the text.} Every figure and table is made by a script listed in \texttt{figures/MANIFEST.md} from \texttt{results/summary}; the numbers quoted in the text are read from \texttt{results/} by \texttt{paper/check\_numbers.py}, which fails if a quoted value differs from the committed tables.

\textbf{Access to excluded material.} The task excludes the published paper on this exact question. No literature search on the question itself was run for this report; the bibliographic look-ups (arXiv and Crossref metadata of the starter-list works and standard references for graph generators and datasets) did not encounter it. The simulation step reports no encounter either (\texttt{REPRODUCE.md} \S8), and the literature-review step's brief states that the question was deliberately not searched and records no encounter.

\section{Additional figures}
\label{app:figs}
\begin{figure}[tbp]
\centering\includegraphics[width=\linewidth]{fig1_strength_all.pdf}
\caption{Bias and RMSE against interference strength $\gamma$ on all six graphs (linear truth, Bernoulli(0.5)).}
\label{fig:strength-all}
\end{figure}
\begin{figure}[tbp]
\centering\includegraphics[width=\linewidth]{fig8_factor_ranking.pdf}
\caption{One-factor ranges (not a ranking; see Sec.~\ref{sec:res-fail}). Range of the RMSE of OLS$(z,f)$ (filled) and of the correctly specified regression (open) over each one-factor sweep; ratio of max to min on the right.}
\label{fig:rank}
\end{figure}
\begin{figure}[tbp]
\centering\includegraphics[width=\linewidth]{fig5_cluster_design.pdf}
\caption{Graph (solid) versus size-matched random (dashed) clusters across resolution, concave spillover: share of units with $f\ge0.9$ (exposure support), bias, RMSE; dotted: Bernoulli(0.5).}
\label{fig:cluster}
\end{figure}
\begin{figure}[tbp]
\centering\includegraphics[width=\linewidth]{fig7_failure_map_relrmse_ba.pdf}
\caption{Relative RMSE failure map on BA (same layout as Fig.~\ref{fig:fail}).}
\label{fig:fail-ba}
\end{figure}
\begin{figure}[tbp]
\centering\includegraphics[width=\linewidth]{fig7_failure_map_relrmse_er.pdf}
\caption{Relative RMSE failure map on ER (same layout as Fig.~\ref{fig:fail}).}
\label{fig:fail-er}
\end{figure}
\begin{figure}[tbp]
\centering\includegraphics[width=\linewidth]{fig7_failure_map_relrmse_ws.pdf}
\caption{Relative RMSE failure map on WS (same layout as Fig.~\ref{fig:fail}).}
\label{fig:fail-ws}
\end{figure}
\begin{figure}[tbp]
\centering\includegraphics[width=\linewidth]{fig7_failure_map_relrmse_stanford3.pdf}
\caption{Relative RMSE failure map on Stanford3 (same layout as Fig.~\ref{fig:fail}).}
\label{fig:fail-stanford3}
\end{figure}
\begin{figure}[tbp]
\centering\includegraphics[width=\linewidth]{fig7_failure_map_relrmse_egofacebook.pdf}
\caption{Relative RMSE failure map on ego-Facebook (same layout as Fig.~\ref{fig:fail}).}
\label{fig:fail-egofacebook}
\end{figure}

\end{document}
